\documentclass[11pt,a4paper]{article}

% --- Language & Font Setup ---
\usepackage[utf8]{inputenc}
\usepackage[T1]{fontenc}
\usepackage[ngerman]{babel}
\usepackage[default]{lato}                    % modern, ATS-safe sans-serif font
\usepackage{iftex}
\ifxetex
  \providecommand{\textls}[2][150]{{\addfontfeature{LetterSpace=15}#2}}
  \renewcommand{\textls}[2][150]{{\addfontfeature{LetterSpace=15}#2}}
\else
  \usepackage[letterspace=150]{microtype}
\fi
\usepackage{xcolor}

% --- Design Colors ---
\definecolor{accent}{HTML}{1F3864}      % dark navy
\definecolor{accentlight}{HTML}{4A6FA5} % lighter blue
\definecolor{datecolor}{HTML}{595959}   % neutral gray

% --- Layout & Margins ---
\usepackage[left=2.5cm, right=2.5cm, top=2.5cm, bottom=2.5cm]{geometry}
\usepackage{hyperref}

% --- Paragraph & Line Spacing ---
\setlength{\parindent}{0pt}
\setlength{\parskip}{7pt}
\setlength{\emergencystretch}{2em}
\linespread{1.15}

% --- ATS & Link Setup ---
\pagestyle{empty}
\urlstyle{same}
\hypersetup{
    colorlinks=true,
    linkcolor=accent,
    urlcolor=accent,
    pdfborder={0 0 0}
}


\begin{document}

% --- Absender & Header ---
    {\Large\bfseries\color{accent}Khaled Swaidan}\\[3pt]
    {\color{datecolor}Schwachhausen, 28211 Bremen \textbullet\ \href{mailto:khaled.swaidan.k@gmail.com}{khaled.swaidan.k@gmail.com}}\\[2pt]
{\color{accent}\rule{\linewidth}{0.8pt}}

    \vspace{0.8cm}

% --- Empfänger ---
    {\color{datecolor}\small Accenture\\
    Recruiting Team\\
    Hamburg}

    \vspace{0.8cm}

% --- Datum & Betreff ---
    \hfill {\color{datecolor}\small Bremen, 28. Oktober 2026}

    \vspace{0.5cm}

    {\large\bfseries\color{accent}Bewerbung als (Junior) AI Native Software Engineer -- Ref. R00359893}

    \vspace{0.5cm}

% --- Anschreiben Text ---
    Sehr geehrte Damen und Herren,

    dass Sie einen „AI Native Software Engineer“ suchen, der KI nicht als Add-on betrachtet, sondern als festen Bestandteil des täglichen Workflows, beschreibt exakt meinen Ansatz. Ich nutze Tools wie ChatGPT, Claude und GitHub Copilot täglich für Code-Generierung, Debugging und Testing – aber ich verantworte auch jede Zeile, die ich in die Produktion bringe.

    In meiner Masterarbeit (Note 1,0) habe ich bewiesen, dass ich die Qualität von KI-Outputs systematisch bewerten und absichern kann. Ich habe einen eigenen 30-Fragen-Benchmark entwickelt, um vier verschiedene LLMs zu evaluieren. Zudem habe ich einen Python-basierten AST-Parser geschrieben, der die Ausgaben der Modelle in Echtzeit gegen eine strikte Grammatik prüft (Guardrails), bevor sie den Nutzer erreichen. Diese Fähigkeit, KI-generierten Code zu verstehen, zu prüfen und zu verbessern, ist mein daily business.

    Die Integration von LLM-APIs in produktive Anwendungen ist für mich kein Neuland. Ich habe eine vollständige RAG-Pipeline (LangChain, ChromaDB) gebaut, inklusive Full-Stack-Entwicklung (FastAPI im Backend, Flutter im Frontend) und Docker-basiertem Deployment. Dabei habe ich gelernt, mit Token-Limits, Latenzzeiten und den Grenzen von Modellen umzugehen.

    Die Zusammenarbeit mit Kunden und die Strukturierung von Anforderungen (Requirements Engineering) habe ich als IT-Berater bei der IFIB Consult gelernt. Ich freue mich sehr darauf, diese Kombination aus technischer KI-Exzellenz und Consulting-Kompetenz in Ihre Kundenprojekte einzubringen und mich langfristig in Richtung „Forward Deployed Engineer“ zu entwickeln.

    Ich bin ab sofort verfügbar und freue mich auf ein Gespräch.

    Mit freundlichen Grüßen,

    \vspace{1.2cm}

    Khaled Swaidan

\end{document}